\section{Git: the safety net}
% Block 4 of PLAN.md, 40 min, with a live demo.  Message: git is what makes it
% safe to let an agent touch your files.  About nine commands; the agent may type
% them, you read the diffs.

% ------------------------------------------------------------------------------
\begin{frame}{Why git, and why now}
  \begin{itemize}
    \item An agent edits ten files in ten seconds
    \item Without a record, you cannot see what changed, or go back
    \item<2-> Git: a lab notebook for your code, with an undo button
  \end{itemize}

  \vspace{3ex}
  \onslide<3->{\centering\ckeybox[0.75\linewidth]{Git is what makes it safe to let an agent
    touch your files}\par}
\end{frame}

% ------------------------------------------------------------------------------
\begin{frame}{Git in one picture}{A chain of snapshots}
  \vspace{2ex}
  \begin{center}
    \begin{tikzpicture}[x=33mm]
      \foreach \i/\h/\m in {0/3f2a91c/{Rabi oscillations},1/b71e04d/{Add analytic test},
                           2/9c0d5e2/{Lindblad decay}} {
        \node[commit] (c\i) at (\i,0) {};
        \node[font=\scriptsize\ttfamily,text=unibasGrey,above=1.5mm of c\i] {\h};
        \node[font=\scriptsize,below=1.5mm of c\i] {\m};
      }
      \draw[flow] (c0) -- (c1);
      \draw[flow] (c1) -- (c2);
      \onslide<4->{\node[font=\scriptsize\bfseries,text=unibasPetrol,right=4mm of c2] {main};}
    \end{tikzpicture}
  \end{center}
  \vspace{1ex}

  \begin{itemize}
    \item A \emph{commit}: a snapshot of the whole project, with message, author and date
    \item<2-> Identified by a \emph{hash}, e.g.\ \cmd{9c0d5e2}
    \item<3-> Any snapshot can be inspected or restored at any time
    \item<4-> \emph{Branch}: a name pointing to the latest commit of a line of work
  \end{itemize}
\end{frame}

% ------------------------------------------------------------------------------
\begin{frame}{Three places for a change}
  \vspace{2ex}
  \begin{center}
    \begin{tikzpicture}[node distance=24mm,
        area/.style={box,minimum width=26mm,minimum height=11mm}]
      \node[area] (w) {\textbf{Working folder}\\{\scriptsize what you edit}};
      \node[area,right=of w] (s) {\textbf{Staging area}\\{\scriptsize next snapshot}};
      \node[area=figModel,right=of s] (h) {\textbf{History}\\{\scriptsize commits}};
      \draw[flow] ([yshift=2mm]w.east) -- node[above,font=\scriptsize\ttfamily] {git add} ([yshift=2mm]s.west);
      \draw[flow] ([yshift=2mm]s.east) -- node[above,font=\scriptsize\ttfamily] {git commit} ([yshift=2mm]h.west);
      \onslide<2->{\draw[flow] (h.south) -- ++(0,-7mm) -| node[pos=0.25,below,font=\scriptsize\ttfamily]
        {git restore <file>} (w.south);}
    \end{tikzpicture}
  \end{center}
  \vspace{4ex}

  \begin{itemize}
    \item You choose what goes into each snapshot
    \item<2-> Unwanted edits in the working folder: restore the last snapshot
    \item<3-> \cmd{git status} tells you where everything is; run it often
  \end{itemize}
\end{frame}

% ------------------------------------------------------------------------------
\begin{frame}[fragile]{Your first repository}
\begin{terminal}
$ cd DrivenSpin
$ git init                                  # start recording
$ git status                                # what is new or changed?
$ git add src/DrivenSpin.jl test/runtests.jl
$ git commit -m "Rabi oscillations with analytic test"
$ git log --oneline                         # the history so far
3f2a91c Rabi oscillations with analytic test
\end{terminal}
  \begin{itemize}
    \item<2-> The message says \emph{why}, not only what: you will read it in a year
    \item<3-> Same commands on macOS, Windows and Linux
  \end{itemize}
\end{frame}

% ------------------------------------------------------------------------------
\begin{frame}{Reading a diff}{The single most important skill of this afternoon}
  \vspace{1ex}
  \begin{diffblock}
    \diffhead{\$ git diff}
    \diffhead{--- a/src/DrivenSpin.jl}
    \diffhead{+++ b/src/DrivenSpin.jl}
    \diffhead{@@ -12,6 +12,6 @@ function hamiltonian($\Delta$, $\Omega$)}
    \diffctx{\ \ \ \ $\sigma$x = [0 1; 1 0]}
    \diffctx{\ \ \ \ $\sigma$z = [1 0; 0 -1]}
    \diffdel{\ \ \ \ return $\Delta$/2*$\sigma$z + $\Omega$*$\sigma$x}
    \diffadd{\ \ \ \ return $\Delta$/2*$\sigma$z + $\Omega$/2*$\sigma$x}
    \diffctx{end}
  \end{diffblock}
  \vspace{2ex}

  \begin{itemize}
    \item \textcolor{unibasRed}{\cmd{-}} line removed, \textcolor{unibasPetrol}{\cmd{+}} line
          added; \cmd{@@} says where
    \item<2-> \cmd{git diff}: changes not yet staged; \cmd{git diff --staged}: what will be committed
    \item<3-> VS Code shows the same, side by side (Source Control panel)
  \end{itemize}
\end{frame}

% ------------------------------------------------------------------------------
\begin{frame}{Spot the problem}{A diff an agent might produce}
  \vspace{1ex}
  \begin{diffblock}
    \diffhead{--- a/test/runtests.jl}
    \diffhead{+++ b/test/runtests.jl}
    \diffhead{@@ -14,3 +14,3 @@ @testset "Rabi" begin}
    \diffdel{\ \ \ \ @test excited\_population(H, $\pi$/$\Omega$) $\approx$ 1.0}
    \diffadd{\ \ \ \ @test excited\_population(H, $\pi$/$\Omega$) $\approx$ 1.0 atol=1.0}
    \diffctx{end}
  \end{diffblock}
  \vspace{2ex}

  \begin{itemize}
    \item<2-> \cmd{atol=1.0}: any population between 0 and 2 now passes
    \item<3-> The test is green and checks nothing
    \item<4-> \term{Any change to a test needs a reason you accept}
  \end{itemize}
\end{frame}

% ------------------------------------------------------------------------------
\begin{frame}{Undo}
  \vspace{1ex}
  \renewcommand{\arraystretch}{1.5}
  \begin{tabularx}{\linewidth}{@{}X l@{}}
    \toprule
    \textbf{Situation} & \textbf{Command} \\
    \midrule
    Unwanted edits, not committed          & \cmd{git restore <file>} \\
    \reveal<2>{Added to staging by mistake}  & \reveal<2>{\cmd{git restore --staged <file>}} \\
    \reveal<3>{A bad commit}                 & \reveal<3>{\cmd{git revert <hash>}} \\
    \bottomrule
  \end{tabularx}

  \vspace{2ex}
  \begin{itemize}
    \item<3-> \cmd{revert} adds a new commit that undoes the old one; history is kept
    \item<4-> \cmd{git reset --hard}, \cmd{git push --force}: they destroy work.
          Never let an agent run them without asking
  \end{itemize}
\end{frame}

% ------------------------------------------------------------------------------
\begin{frame}{Branches: cheap experiments}
  \vspace{2ex}
  \begin{center}
    \begin{tikzpicture}[x=24mm,y=11mm]
      \node[commit] (a) at (0,0) {};
      \node[commit] (b) at (1,0) {};
      \node[commit] (c) at (2,0) {};
      \node[font=\scriptsize\bfseries,text=unibasPetrol,left=3mm of a] {main};
      \draw[flow] (a) -- (b); \draw[flow] (b) -- (c);
      \onslide<2->{%
        \node[branchcommit] (d) at (2,1) {};
        \node[branchcommit] (e) at (3,1) {};
        \draw[flow,draw=unibasTurquoise] (b) -- (d); \draw[flow,draw=unibasTurquoise] (d) -- (e);
        \node[font=\scriptsize\bfseries,text=unibasTurquoise,above=2mm of e] {try-rk4};}
      \onslide<3->{%
        \node[commit] (f) at (4,0) {};
        \draw[flow] (c) -- (f); \draw[flow,draw=unibasTurquoise] (e) -- (f);
        \node[font=\scriptsize,below=1.5mm of f] {merge};}
    \end{tikzpicture}
  \end{center}
  \vspace{1ex}

  \begin{itemize}
    \item<2-> \cmd{git switch -c try-rk4}: a new line of work; \cmd{main} stays untouched
    \item<3-> It worked: \cmd{git switch main}, then \cmd{git merge try-rk4}
    \item<4-> It did not: \cmd{git branch -D try-rk4}, and it never happened
  \end{itemize}
\end{frame}

% ------------------------------------------------------------------------------
\begin{frame}{Remotes: GitHub}
  \vspace{2ex}
  \begin{center}
    \begin{tikzpicture}[node distance=22mm]
      \node[box] (me) {\textbf{Your laptop}\\{\scriptsize full history}};
      \node[box=figModel,right=of me] (gh) {\textbf{GitHub}\\{\scriptsize full history}};
      \node[box,right=of gh] (col) {\textbf{Colleague}\\{\scriptsize full history}};
      \draw[flow] ([yshift=2mm]me.east) -- node[above,font=\scriptsize\ttfamily] {git push} ([yshift=2mm]gh.west);
      \draw[flow] ([yshift=-2mm]gh.west) -- node[below,font=\scriptsize\ttfamily] {git pull} ([yshift=-2mm]me.east);
      \draw[flow] (gh) -- node[above,font=\scriptsize\ttfamily] {git clone} (col);
    \end{tikzpicture}
  \end{center}
  \vspace{2ex}

  \begin{itemize}
    \item Backup, sharing, and one place for the group's code
    \item<2-> Login over SSH, set up in the installation guide
    \item<3-> GitHub can run your tests on every push: after the break
  \end{itemize}
\end{frame}

% ------------------------------------------------------------------------------
\begin{frame}[fragile]{What goes into git}
  \begin{columns}[T]
    \begin{column}{0.46\textwidth}
      \textbf{Commit}
      \begin{itemize}
        \item Source, scripts, tests
        \item \cmd{Project.toml}, \cmd{Manifest.toml}
        \item \cmd{README.md}, \cmd{AGENTS.md}
      \end{itemize}
    \end{column}
    \begin{column}{0.46\textwidth}
      \textbf{Do not commit}
      \begin{itemize}
        \item Large raw data: archive it elsewhere
        \item Generated plots and results
        \item Passwords and API keys, \emph{ever}
      \end{itemize}
    \end{column}
  \end{columns}
  \vspace{1ex}
\begin{onlyenv}<2->
\begin{terminal}[title={\hfill\tiny\color{unibasGrey}\smallcapsline{.gitignore}}]
data/
plots/
*.jld2
\end{terminal}
\end{onlyenv}
\end{frame}

% ------------------------------------------------------------------------------
\begin{frame}[fragile]{Git with an agent}{Habits}
  \begin{enumerate}
    \item Commit before the agent starts: a clean point to return to
    \item<2-> One task, one commit
    \item<3-> Read the diff before every commit, tests first
    \item<4-> Let the agent type the commands; you check the result
  \end{enumerate}
\end{frame}

% ------------------------------------------------------------------------------
\begin{frame}[fragile]{Let the agent keep the record}{Prompts that use git}
\begin{prompt}
Commit my manual changes before starting.
\end{prompt}
\begin{visibleenv}<2->
\begin{prompt}
Work in modular commits: one commit per change, with the tests passing at each.
\end{prompt}
\end{visibleenv}
\begin{visibleenv}<3->
\begin{terminal}
$ git log --oneline
e4b1a20 Add Ramsey sequence with fringe-frequency test
7c93f01 Add Lindblad dissipator with trace and positivity tests
51d0b8e Relabel plot axes in scripts/rabi.jl (manual changes)
\end{terminal}
\end{visibleenv}
  \begin{itemize}
    \item Your edits and the agent's never end up in the same commit
    \item<2-> Each change can be reviewed, and reverted, on its own
    \item<4-> You still read every one: \cmd{git show <hash>}
  \end{itemize}
\end{frame}

% ------------------------------------------------------------------------------
\begin{frame}[fragile]{The agent wrecks the project}{Live demo}
\begin{prompt}
Refactor the whole project so that all parameters live in one struct.
\end{prompt}
\begin{onlyenv}<2->
\begin{terminal}
$ git status --short          # what did it touch?
 M scripts/rabi.jl
 M src/DrivenSpin.jl
 M test/runtests.jl
$ git diff --stat             # how much?
$ git restore .               # back to the last commit
\end{terminal}
\end{onlyenv}
  \begin{itemize}
    \item<3-> Thirty seconds, nothing lost
    \item<4-> \cmd{git restore .} discards \emph{all} uncommitted changes: check first
  \end{itemize}
\end{frame}

% ------------------------------------------------------------------------------
\begin{frame}{The commands you need}
  \vspace{1ex}
  \renewcommand{\arraystretch}{1.25}
  \begin{tabularx}{\linewidth}{@{}l X@{}}
    \toprule
    \cmd{git clone <address>}   & copy a repository from GitHub \\
    \cmd{git status}            & what changed? \\
    \cmd{git diff}              & how exactly? \\
    \cmd{git add <file>}        & put a change into the next snapshot \\
    \cmd{git commit -m "why"}   & take the snapshot \\
    \cmd{git log --oneline}     & the history \\
    \cmd{git restore <file>}    & throw away uncommitted edits \\
    \cmd{git switch -c <name>}  & start a branch \\
    \cmd{git push} / \cmd{git pull} & sync with GitHub \\
    \bottomrule
  \end{tabularx}
\end{frame}
